\chapter{Téléopération de Protova2 : WiFi et 5G}
\chapminitoc
\label{chap:teleop}

\section{Contexte}

Avant de rendre un véhicule autonome, il faut pouvoir le conduire à distance :
la téléopération est à la fois un mode d'exploitation à part entière (conduite
déportée avec retour vidéo), un outil de développement indispensable
(cartographie manuelle pour le SLAM, essais de perception) et le mode de repli
de sécurité lorsque l'autonomie est en défaut. Sur Protova2, le poste de
pilotage est un PC équipé d'un volant à retour de force Logitech G29 ; le
véhicule embarque la Jetson Orin Nano, qui pilote les actionneurs via le
microcontrôleur Pico et renvoie le flux de sa caméra.

J'ai mis en œuvre et comparé deux liaisons radio pour cette téléopération : le
\textbf{WiFi} du laboratoire, puis un \textbf{réseau 5G privé}, représentatif
des déploiements industriels visés pour les véhicules connectés.

\section{Problématique}

La téléopération est une \emph{boucle} : l'opérateur perçoit la scène à
travers la vidéo, agit sur le volant, et sa commande traverse le réseau avant
d'atteindre les roues. Ce que voit l'opérateur est donc toujours \emph{le
passé} du véhicule, et sa commande arrive toujours \emph{en retard} sur ce
qu'il a vu. La qualité de conduite dépend ainsi de trois grandeurs : la
\textbf{latence} (le retard de la boucle), le \textbf{débit} (la capacité à
transporter la vidéo) et la \textbf{cadence d'images} (la fluidité du retour
visuel). L'un des enjeux du chapitre est précisément de déterminer laquelle de
ces grandeurs est réellement dimensionnante.

À cette problématique de performance s'ajoute une problématique
d'intégration : le middleware de ROS~2 (DDS) repose sur la découverte
\emph{multicast}, qui fonctionne nativement sur un réseau local WiFi mais
\textbf{ne traverse pas un réseau cellulaire}. La liaison 5G impose donc une
architecture de transport différente de celle du WiFi.

\section{Architecture logicielle de commande}

La chaîne de commande est identique quel que soit le support radio ; seul le
transport change. Côté PC, \texttt{joy\_node} lit le volant et publie
\texttt{/joy} ; côté véhicule, \texttt{g29\_teleop} convertit ces axes en
consigne de conduite (\texttt{/cmd\_vel\_G29}), un multiplexeur
(\texttt{twist\_mux}) arbitre entre les sources de commande par priorité ---
la téléopération, la navigation autonome et le freinage d'urgence (priorité
maximale, chapitre~\ref{chap:perception}) --- et le nœud série \texttt{TX}
transmet la consigne retenue au Pico, qui commande le servomoteur de direction
et le variateur de traction. En sens inverse, la caméra Orbbec publie son flux
couleur compressé JPEG, visualisé sur le PC. Chaque entrée du multiplexeur est
assortie d'un \emph{timeout} : une source qui cesse d'émettre perd la main,
ce qui garantit qu'une coupure réseau conduit à l'arrêt et non à la poursuite
de la dernière commande reçue.

\begin{figure}[ht]
  \centering
  \begin{tikzpicture}[x=1mm, y=1mm,
      boite/.style={draw=altenink, fill=altenink!8, rounded corners=2pt,
                    minimum height=8mm, minimum width=22mm, align=center, font=\small},
      lien/.style={-{Stealth[length=2.5mm]}, thick, color=black!55},
      etiq/.style={font=\tiny\ttfamily, color=black!55}]
    \node[boite] (joy)  at (0,0)    {volant G29\\+ \texttt{joy\_node} (PC)};
    \node[boite] (g29)  at (40,0)   {\texttt{g29\_teleop}};
    \node[boite] (mux)  at (76,0)   {\texttt{twist\_mux}\\(priorités)};
    \node[boite] (tx)   at (110,0)  {\texttt{TX} série};
    \node[boite] (pico) at (141,0)  {Pico\\servo + ESC};
    \node[boite] (cam)  at (76,-22) {caméra Orbbec\\(JPEG 720p)};
    \node[boite] (pcv)  at (20,-22) {visualisation\\PC};
    \draw[lien] (joy.east)  -- node[etiq, above]{/joy \ (réseau)} (g29.west);
    \draw[lien] (g29.east)  -- node[etiq, above]{/cmd\_vel\_G29} (mux.west);
    \draw[lien] (mux.east)  -- node[etiq, above]{/cmd\_vel} (tx.west);
    \draw[lien] (tx.east)   -- node[etiq, above]{série} (pico.west);
    \draw[lien] (cam.west)  -- node[etiq, above]{images \ (réseau)} (pcv.east);
  \end{tikzpicture}
  \caption{Chaîne de téléopération : les commandes descendent du PC vers le
  véhicule, la vidéo remonte. Les liens marqués « réseau » traversent le WiFi
  ou la 5G selon la configuration.}
  \label{fig:chaine-teleop}
\end{figure}

\section{Mission 1 : Téléopération WiFi (FastDDS)}

\objectif
Conduire Protova2 au volant G29 à travers le WiFi du laboratoire, avec retour
vidéo temps réel, en utilisant le transport natif de ROS~2.

\methodes
Sur un réseau local, le DDS natif de ROS~2 (FastDDS) suffit : les deux
machines partagent le même domaine (\texttt{ROS\_DOMAIN\_ID=125}) et la
découverte multicast fait le reste, sans configuration supplémentaire. J'ai
encapsulé le lancement dans des scripts symétriques
(\texttt{launch\_pc\_wifi.sh} côté poste, \texttt{launch\_car\_wifi.sh} côté
véhicule) qui fixent l'environnement, nettoient les instances précédentes et
démarrent les nœuds dans l'ordre. Un point de vigilance : les flux d'images
doivent être publiés et souscrits en QoS \emph{best-effort} --- en mode
\emph{reliable}, les retransmissions DDS s'accumulent sur un lien radio et le
flux se fige.

\resultats
La conduite au volant fonctionne avec un retour vidéo fluide. La
caractérisation du lien (protocole détaillé en
section~\ref{sec:campagne-mesures}) donne : latence moyenne de \textbf{68~ms}
mais très irrégulière (pointes à 237~ms), débit d'environ \textbf{185~Mbps}
dans les deux sens, et cadence vidéo de \textbf{30~images/s} --- la pleine
cadence de la caméra.

\discussion
Le débit WiFi est très confortable, mais la latence présente une particularité
contre-intuitive que mes mesures ont mise en évidence : elle est \emph{pire à
vide qu'en charge} (84~ms au repos contre 59~ms sous trafic soutenu). Il
s'agit de l'économie d'énergie de la radio WiFi : entre deux paquets, elle
s'endort, et son réveil coûte 100 à 300~ms. Ce comportement, sans conséquence
pour de la bureautique, est pénalisant pour une boucle de pilotage, où les
commandes sont justement de petits paquets épars.

\section{Mission 2 : Téléopération 5G (Zenoh)}

\objectif
Reproduire la téléopération complète à travers un réseau 5G privé, en
remplaçant le transport DDS, inopérant sur un lien cellulaire.

\methodes
Chaque machine est équipée d'un modem 5G Quectel (RM520N côté PC, RM530N côté
véhicule). Le multicast étant impossible sur ce lien, j'ai remplacé le
transport par \textbf{Zenoh} (\texttt{rmw\_zenoh\_cpp}) en architecture
\emph{hub-and-spoke} : un routeur Zenoh sur le PC écoute sur le port 7447 ; le
routeur du véhicule ouvre une connexion TCP vers lui à travers la 5G ; et les
nœuds ROS~2 de chaque machine se connectent en clients à leur routeur local.
Deux subtilités de configuration méritent d'être notées : \texttt{rmw\_zenoh}
ignore la variable \texttt{ZENOH\_CONFIG} et exige
\texttt{ZENOH\_ROUTER\_CONFIG\_URI} / \texttt{ZENOH\_SESSION\_CONFIG\_URI} ;
et les sessions des nœuds ne se rattachent pas à un routeur redémarré, ce qui
impose un ordre de démarrage strict.

\resultats
La chaîne complète fonctionne à travers la 5G : appairage des routeurs vérifié
(session TCP établie sur le port 7447 à travers le lien radio), commandes du
volant reçues par le véhicule, vidéo remontée au poste. Les performances
mesurées figurent en section~\ref{sec:campagne-mesures}.

\discussion
Le débogage le plus instructif a concerné l'\textbf{adressage IP}. Les
symptômes : le lien 5G semblait mort --- \texttt{ping 172.16.48.7} répondait
\emph{« Do you want to ping broadcast? »} --- alors que la radio était
parfaitement fonctionnelle. Le diagnostic : le modem attribue par DHCP une
adresse en \texttt{/30}, un sous-réseau de quatre adresses seulement, dans
lequel \texttt{.7} --- précisément l'adresse du poste de pilotage --- est
l'\emph{adresse de diffusion} : le noyau refuse alors tout trafic unicast vers
le PC. La correction consiste à élargir le masque en \texttt{/28}
(\texttt{.7} redevient une adresse d'hôte ordinaire) et à purger
systématiquement la \texttt{/30} que chaque renouvellement de bail DHCP
réintroduit.

Plus généralement, aucune des difficultés rencontrées n'incombait au protocole
Zenoh lui-même : toutes relevaient de l'\emph{intégration système} ---
installation, variables d'environnement, adressage, ordonnancement du
démarrage, et jusqu'à ce piège d'exploitation : un \texttt{pkill}
lancé à travers SSH avec le nom d'un nœud en argument peut tuer sa propre
session, le motif se retrouvant dans la ligne de commande du shell distant.
La méthode qui a permis d'en venir à bout est la vérification \emph{maillon
par maillon} : interface $\rightarrow$ adresse IP $\rightarrow$ \texttt{ping}
$\rightarrow$ session TCP 7447 $\rightarrow$ graphe ROS. Enfin, pour rendre
l'ensemble robuste aux aléas, j'ai confié la pile du véhicule à un service
\textbf{systemd} (\texttt{car5g.service}) : elle démarre au boot, se relance
seule en cas d'échec et ne dépend plus d'aucun terminal ouvert.

\section{Campagne de mesures : WiFi vs 5G}
\label{sec:campagne-mesures}

\objectif
Comparer objectivement les deux liaisons sur les trois grandeurs de la
téléopération --- latence, débit, cadence vidéo --- avec des outils
reproductibles.

\methodes
J'ai développé trois scripts de mesure autonomes, chacun produisant sa courbe
et ses données brutes : \texttt{mesure\_latence.py} (RTT de chaque paquet
\texttt{ping}, 5~Hz), \texttt{mesure\_debit.py} (débit utile transféré chaque
seconde par \texttt{iperf3}, dans les deux sens) et \texttt{mesure\_fps.py}
(comptage, seconde par seconde, des images reçues au poste de pilotage). Les
mêmes scripts servent pour les deux réseaux --- seule change la cible --- ce
qui garantit une comparaison à protocole identique. Ils complètent les
tableaux de bord de supervision en direct présentés au
chapitre~\ref{chap:dashboard}.

\resultats
La figure~\ref{fig:comparaison-wifi-5g} superpose les mesures ; le
tableau~\ref{tab:wifi-5g} les résume.

\begin{figure}[ht]
  \centering
  \includegraphics[width=\textwidth]{comparaison_wifi_5g.png}
  \caption{Comparaison WiFi / 5G sur le lien PC $\leftrightarrow$ véhicule :
  latence (gauche), débit dans les deux sens (centre), cadence vidéo
  (droite).}
  \label{fig:comparaison-wifi-5g}
\end{figure}

\begin{table}[ht]
  \centering
  \begin{tabular}{lccc}
    \hline
    \textbf{Grandeur} & \textbf{WiFi} & \textbf{5G} & \textbf{Écart} \\
    \hline
    Latence moyenne (gigue)          & 68 ms (pointes 237)  & \textbf{26 ms} (17--37) & 5G $\times$2,6 \\
    Débit montant PC$\to$véhicule    & \textbf{186 Mbps}    & 69 Mbps   & WiFi $\times$2,7 \\
    Débit descendant véhicule$\to$PC & \textbf{183 Mbps}    & 77 Mbps   & WiFi $\times$2,4 \\
    Cadence vidéo 720p               & 30,5 im/s            & 30,5 im/s & égalité \\
    \hline
  \end{tabular}
  \caption{Synthèse des mesures WiFi / 5G (moyennes des campagnes).}
  \label{tab:wifi-5g}
\end{table}

\discussion
Ces résultats illustrent une distinction essentielle : le débit et la cadence
d'images sont des critères \emph{à seuil}, tandis que la latence est le
\emph{facteur de mérite}. Le besoin réel du système est modeste --- environ
18~Mbps pour la vidéo 720p, moins de 0,1~Mbps pour les commandes --- et les
deux liens le dépassent largement : la vidéo passe d'ailleurs à pleine cadence
(30~im/s) sur les deux. L'avantage de débit du WiFi ne se traduit donc par
aucun bénéfice perceptible. La latence, en revanche, intervient dans la boucle
perception--action à chaque instant et ne peut être compensée par aucun autre
paramètre : à 1~m/s, les 68~ms moyens du WiFi (et surtout ses pointes à
237~ms) représentent 7 à 24~cm parcourus « à l'insu » de l'opérateur, contre
moins de 3~cm --- et de façon \emph{prévisible} --- en 5G. Pour la
téléopération temps réel, la 5G est donc le lien le plus adapté, par sa
latence faible et \emph{stable} bien plus que par toute autre caractéristique.

Une perspective naturelle est la \textbf{mobilité} : le \emph{handover}
(changement de point d'accès en mouvement) est natif et quasi transparent en
5G --- c'est l'un de ses atouts fondateurs pour le véhicule connecté ---
tandis que le \emph{roaming} WiFi, décidé par le client, provoque des coupures
de quelques centaines de millisecondes à plusieurs secondes. Le bâtiment du
laboratoire, couvert par une vingtaine de bornes WiFi, se prête à une mesure
comparative de ces interruptions, qui constituerait un prolongement direct de
cette campagne.

\section{Conclusion}

La téléopération de Protova2 fonctionne de bout en bout sur les deux liaisons,
avec une chaîne de commande commune et deux transports adaptés à chaque
support (FastDDS en WiFi, Zenoh en hub-and-spoke sur 5G). La campagne de
mesures, menée avec des outils développés pour être reproductibles, dégage la
conclusion structurante du chapitre : à débit suffisant --- ce qui est le cas
des deux liens --- c'est la latence, et sa stabilité, qui gouverne la qualité
de la téléopération ; la 5G y est nettement supérieure (26~ms stables contre
68~ms irréguliers). Ce résultat, ainsi que la robustesse apportée par la
supervision systemd de la pile embarquée, préparent le terrain du
chapitre~\ref{chap:v2v}, où le lien inter-véhicules transporte non plus des
commandes de conduite mais des messages de sécurité coopérative.
